\section{Introduction}
The $n$-queens problem asks in how many ways $n$ queens can be placed on an $n\times n$ board so
that no two attack one another. We denote this number by $Q(n)$ and count solutions related by
symmetry separately. Reading the column of each queen in row order identifies each solution with a
permutation. The diagonal constraints determine which of the $n!$ permutations qualify.

Simkin~\cite[Theorem~1.1]{simkin} obtained the leading order of $\log Q(n)$ from a variational
problem. Luria~\cite{luria} obtained upper bounds by the entropy method. Luria and
Simkin~\cite{ls} and, independently, Bowtell and Keevash~\cite{bk} proved the lower bound
$Q(n)\ge((1-o(1))ne^{-3})^n$, which confirms a conjecture of Rivin, Vardi and
Zimmermann~\cite{rvz}. Nobel, Agrawal and Boyd~\cite{nab} bounded the first-order constant.
Liu, Liao and Wang~\cite{llw} studied the problem as a lattice gas and estimated that constant
numerically from Monte Carlo simulations. The
completion of partial configurations, which arises here in every residual board, is studied
in~\cite{gms}, and the history of the problem is surveyed in~\cite{bs}. A sharper estimate is
needed to determine the multiplicative factor that remains in the count after the leading term is
extracted.

Let $Q_{\rm pred}(n)$ be the canonical predictor of~\cite{pred} and write
$R_Q(n)=\log Q(n)-\log Q_{\rm pred}(n)$. The following result establishes that $R_Q(n)\to0$, the
qualitative form of Conjecture~10.1 of~\cite{pred}.
\begin{teorema}[Constant term of the count]\label{princ}
On the ordinary board, as $n\to\infty$,
\[
 \log Q(n)-\log n!-n\Phi\longrightarrow
 \tfrac12\log2-\tfrac12-\tfrac32 I_2-\tfrac12\log{\det}_2(I+J)=C_{\rm pred}.
\]
\end{teorema}
Here $\nu_1$ is the optimal profile of mass one from~\cite{perfil}, $\Phi=s_1+1$ with $s_1$ its
entropy, and $I_2=\int_{[0,1]^2}\nu_1^2$. The symbol ${\det}_2$ denotes the regularized
Carleman--Fredholm determinant, ${\det}_2(I+J)=\det\bigl((I+J)e^{-J}\bigr)$, defined for every
Hilbert--Schmidt operator $J$ \cite[Chapter~9]{simon}; this is the convention
of~\cite[(D4.1)]{pred}. The operator $J$ of the theorem is described below.
Put $\alpha=-s_1$, so $\Phi=1-\alpha$. Stirling's formula
turns the theorem into an asymptotic formula refining Simkin's leading-order estimate:
\[
 Q(n)\sim e^{C_{\rm pred}}\,n!\,e^{n\Phi}
      \sim e^{C_{\rm pred}}\sqrt{2\pi n}\,(n e^{-\alpha})^n .
\]
Since \cite[Theorem~1.1]{simkin} gives $Q(n)^{1/n}/n\to e^{-\alpha_{\rm S}}$ for a constant
$\alpha_{\rm S}$, the theorem implies $\alpha=\alpha_{\rm S}$: the number $-s_1$ is Simkin's
constant, which the bounds of~\cite{nab} place in $[1.944000752,\,1.944001082]$.
Subtracting $\log n!+n\Phi$ from $\log Q(n)$ leaves the constant $C_{\rm pred}$ in the limit;
subtracting $n\log n-\alpha n+\tfrac12\log n$ leaves $C_{\rm pred}+\tfrac12\log(2\pi)$.
The theorem gives $Q(n)/Q_{\rm pred}(n)\to1$. The quantitative estimate $|R_Q(n)|\le C/n$
in the same conjecture remains open.

The operator $J=\mathcal K+P_0$ acts on two parity sheets for each line family (rows, columns,
diagonals and antidiagonals), with measure $dt/2$ on each sheet. The off-diagonal kernels of
$\mathcal K$ evaluate $\nu_1$ at the intersection of the two lines when it lies on the board.
The diagonal--antidiagonal block couples only opposite sheets: the indices $i-j$ and $i+j-1$
sum to the odd integer $2i-1$. The other five blocks are independent of the sheet indices.
The projection $P_0$ is onto the unit vector taking the values $1/\sqrt2$ on the row sheets,
$-1/\sqrt2$ on the column sheets, and zero on the other sheets
\cite[\S5.3, (D3.4)--(D3.6)]{pred}. The uniform gap gives $I+J\ge aI$ for some $a>0$;
the regularized logarithm in the theorem is therefore real \cite[\S5.4, (D4.1)]{pred}.

Section~9 of~\cite{pred} gives the estimate $C_{\rm pred}\approx0.9401$, and direct evaluation of
the formula gives $0.940144$. This evaluation \cite[Section~9, ``Direct evaluation of the
formula'']{pred} uses a Nyström discretization of the operator of Section~5.3 of that work, on
successive grids and with extrapolation in the step size; the pairs of grids of sizes 192 and 384,
and 193 and 385, gave $0.94014395698$ and $0.94014395711$. None of these values has a certified
error bound. With the reference value
$\Phi_0=1-1.944001$ from that section, the exact counts for $16\le n\le27$ \cite{oeis} give
values of $\log Q(n)-\log n!-n\Phi_0$ between $0.913$ and $0.941$ \cite[Table~2]{pred}; the value
$Q(27)$ comes from~\cite{q27} and was verified independently in~\cite{yaoli}. In that range, when
$C_{\rm pred}$ is replaced by $0.940144$, the ratio $Q(n)/(e^{C_{\rm pred}}n!\,e^{n\Phi_0})$ lies
between $0.973$ and $1.001$. The residual computed with $\Phi_0$ minus the one computed with
$\Phi$ is exactly $n(\Phi-\Phi_0)$. Using numerical
values of the predictor, the same section gives $R_Q(20)\approx0.0208$ and $R_Q(27)\approx0.0119$.
These values are consistent with Theorem~\ref{princ}, which establishes a limit without giving a
rate of convergence.

Section~10 of~\cite{pred} identifies three questions. The first concerns the regularity of the
residual boards along the random process. The second concerns the effect of the last few rows, when
few insertion choices remain. The third is how to combine the errors from successive insertions
when comparing the two counts.

\paragraph{Outline of the proof.}
Write $g(B)=\log Q(B)-L(B)$, where $L(B)$ is the predictor for a residual board $B$
(Section~\ref{leyes}). We must prove that $g(B_0)\to0$ on the initial board $B_0$. We compare two
laws on histories that insert one queen per row in a uniform row order. The uniform law
$\mu$ reveals a uniform solution. The adaptive law $\nu$ chooses each square $e$ with probability
proportional to $\exp L(B+e)$, the prediction for the resulting board. Section~\ref{historias}
expresses $g(B_0)$ in terms of the two laws on a stopped history: an identity involving the relative
entropy of $\mu$ with respect to $\nu$ and a lower bound under $\nu$. The accumulated term in these
formulas is small by the consistency result of~\cite{pred}.

Fix a small $\delta>0$ and a large integer $K$, and set $M=\lceil\delta n\rceil$.
Section~\ref{meso:constantes} specifies the order in which the constants are chosen.
Sections~\ref{macro} and~\ref{meso} prove that, under $\nu$, the boards in the history are regular
(Definition~\ref{def:regular}) until $K$ empty rows remain, except with probability
$O(e^{-cK})+o(1)$. Down to $M$ empty rows we track the availability of lines and squares
against a deterministic trajectory. From $M$ to $K$ we control diagonal mass and preserve the
diffusion, vacancy, and gap bounds needed for consistency.
Section~\ref{costes} bounds $g$ from above on the first board that ceases to be regular.
Write $B^{(K)}$ for the board with $K$ empty rows and $F$ for the event that every board up to
and including $B^{(K)}$ is regular. We call the intersection of an empty row and an empty column an
empty crossing. Let $D_K$ be the event that two empty crossings of
$B^{(K)}$ share a diagonal. Section~\ref{terminal} proves that
$\nu(F\cap D_K)=O(K^4/n)$ and that, under $\mu$, the uniform row order gives
$\mu(D_K\mid\pi)=O(K^4/n)$ for every fixed solution $\pi$. On $F\cap D_K^c$ the count is a
permanent and $|g(B^{(K)})|\le\theta_K$, with $\theta_K\to0$ as $K\to\infty$. In Section~\ref{balance}
we first let $n\to\infty$ with $\delta$ and $K$ fixed, and then let $K\to\infty$.
Figure~\ref{fig:tramos} summarizes the three ranges.

Theorems~1.1--1.3, Proposition~5.2 and Theorem~7.7 of~\cite{perfil} give the continuous profile
and its bounds. The entropy comparison in Theorem~3.6 of~\cite{ultimas} motivates the argument in
Section~\ref{costes}, where we verify the estimates on the support of each residual board. We use
Theorems~7.1 and~8.1 and Corollary~7.2 of~\cite{pred} for insertions, and Theorems~1.1 and~5.7 of
that work for the constant. The only asymptotic estimate for permanents that we use is that of
McCullagh~\cite{mccullagh}, in the uniform form proved by Divoux and
Wang~\cite[Proposition~3.1 and Remark~2.1]{dw}.
In Section~\ref{permanente} we verify its hypotheses after scaling the matrix of the residual
board.
